\begin{examplebox}
	\textbf{\textcolor{CExample}{例1（基础）}}

	\textbf{\textcolor{CExample}{题目：}}
	已知正四面体的棱长为$6$，求它的高。

	\textbf{\textcolor{CExample}{解题步骤：}}
	\begin{description}[style=nextline, leftmargin=2.8em, labelsep=0.5em, itemsep=0.45em]
		\item[\textbf{第一步（代公式）：}]写出正四面体高公式$h = \frac{\sqrt{6}}{3} a$.
			\par\textbf{理由：}此为正四面体高公式，属于已知结论。
		\item[\textbf{第二步（代入数值）：}]将$a=6$代入：$h = \frac{\sqrt{6}}{3} \times 6 = 2\sqrt{6}$.
			\par\textbf{理由：}直接数值计算，分母$3$与$6$约分得$2$.
		\item[\textbf{第三步（得结果）：}]所以该正四面体的高为$2\sqrt{6}$.
			\par\textbf{理由：}化简完成，结果已得到。
	\end{description}

	\textbf{\textcolor{CExample}{关键结论：}}
	\textbf{$h = 2\sqrt{6}$}.

	\medskip
	\textbf{\textcolor{CExample}{例2（稍变形）}}

	\textbf{\textcolor{CExample}{题目：}}
	已知正四面体的高为$\sqrt{6}$，求它的体积。

	\textbf{\textcolor{CExample}{解题步骤：}}
	\begin{description}[style=nextline, leftmargin=2.8em, labelsep=0.5em, itemsep=0.45em]
		\item[\textbf{第一步（反求棱长）：}]由$h = \frac{\sqrt{6}}{3} a$反解得$a = \frac{3h}{\sqrt{6}}$. 代入$h=\sqrt{6}$得$a = \frac{3\sqrt{6}}{\sqrt{6}} = 3$.
			\par\textbf{理由：}高公式变形，消去$\sqrt{6}$.
		\item[\textbf{第二步（写体积公式）：}]正四面体体积公式$V = \frac{\sqrt{2}}{12} a^3$.
			\par\textbf{理由：}体积公式可由高与底面积推导得到，是常见变形结论.
		\item[\textbf{第三步（代入得体积）：}]将$a=3$代入：$V = \frac{\sqrt{2}}{12} \times 27 = \frac{27\sqrt{2}}{12} = \frac{9\sqrt{2}}{4}$.
			\par\textbf{理由：}分数约简即可.
	\end{description}

	\textbf{\textcolor{CExample}{关键结论：}}
	\textbf{$V = \frac{9\sqrt{2}}{4}$}.

	\medskip
	\textbf{\textcolor{CExample}{例3（综合应用）}}

	\textbf{\textcolor{CExample}{题目：}}
	已知正四面体的表面积为$36\sqrt{3}$，求它的高。

	\textbf{\textcolor{CExample}{解题步骤：}}
	\begin{description}[style=nextline, leftmargin=2.8em, labelsep=0.5em, itemsep=0.45em]
		\item[\textbf{第一步（由表面积求棱长）：}]正四面体表面积$S_{\text{表}} = 4 \times \frac{\sqrt{3}}{4} a^2 = \sqrt{3} a^2$. 由已知$S_{\text{表}} = 36\sqrt{3}$得$\sqrt{3} a^2 = 36\sqrt{3}$，所以$a^2 = 36$，$a = 6$.
			\par\textbf{理由：}等边三角形面积公式乘以$4$，解方程即可.
		\item[\textbf{第二步（代高公式）：}]将$a=6$代入$h = \frac{\sqrt{6}}{3} a$得$h = \frac{\sqrt{6}}{3} \times 6 = 2\sqrt{6}$.
			\par\textbf{理由：}直接使用正四面体高公式.
		\item[\textbf{第三步（得结果）：}]所以所求高为$2\sqrt{6}$.
			\par\textbf{理由：}计算完毕，答案已得.
	\end{description}

	\textbf{\textcolor{CExample}{关键结论：}}
	\textbf{$h = 2\sqrt{6}$}.
\end{examplebox}
